\documentclass[11pt]{article}
\usepackage[a4paper,margin=26mm]{geometry}
\usepackage{amsmath,amssymb,array}
\setlength{\parindent}{0pt}
\setlength{\parskip}{7pt}
\emergencystretch=2em
\title{Three distinct voters in a one-dimensional\\single-peaked, single-crossing profile\\
\large Disproof of conjecture 00000008244}
\author{ziangni-sys}
\date{4 October 2026}
\begin{document}
\maketitle
The source's middle conjunct states that the maximal agent count
in the simultaneous single-peaked and single-crossing domain is
$2^d d$. In dimension one this is two.
An actual one-dimensional Euclidean profile with three distinct
strict preferences disproves that bound. We use neither duplicate
preferences nor a higher-dimensional embedding.

\textbf{Actual locations and utilities.}
There are three alternatives $a,b,c$ with locations
\[
 q_a=0,\qquad q_b=2,\qquad q_c=3
\]
on the real line. Three agents have distinct locations
\[
 p_0=0,\qquad p_1=5/4,\qquad p_2=7/4.
\]
The actual Euclidean utility of agent $i$ for alternative $z$ is
\[
 u_i(z)=-(q_z-p_i)^2.
\]
These give the complete table
\[
\begin{array}{c|rrr|l}
 i & u_i(a)&u_i(b)&u_i(c)&\text{strict order}\\ \hline
 0 & 0 & -4 & -9 & a\succ b\succ c\\
 1 & -25/16 & -9/16 & -49/16 & b\succ a\succ c\\
 2 & -49/16 & -1/16 & -25/16 & b\succ c\succ a
\end{array}
\]
Every agent has a strict total order, with no ties, and all three
orders are distinct. The alternative axis is $a<b<c$; the voter
order is $0<1<2$, agreeing with their actual real locations.

\textbf{Single-peakedness.}
Agent zero has peak $a$ and prefers alternatives less as one moves
right. Agents one and two have peak $b$, with $b$ preferred to
the sole alternative on each side. Equivalently, each preference
increases toward its peak along the fixed axis and decreases away
from it on either side. Thus the complete profile is single-peaked
on this one-dimensional axis.

\newpage
\textbf{Single-crossing, including all alternative pairs.}
Along the voter order $0<1<2$, the strict-comparison sets are
\[
 \{i:a\succ_i b\}=\{0\},\qquad
 \{i:a\succ_i c\}=\{0,1\},\qquad
 \{i:b\succ_i c\}=\{0,1,2\}.
\]
Each is an initial segment. Each reversed comparison is its
complement, a final segment; a comparison of an alternative to
itself is empty. Hence every ordered-pair comparison set is an
initial or final segment of the voter order. This is the standard
single-crossing condition: a pairwise preference can switch at
most once. In particular $a$ versus $b$ switches after voter zero,
$a$ versus $c$ after voter one, and $b$ versus $c$ never switches.

\textbf{The actual dimension and failure of the count bound.}
Both agents and alternatives are located in $\mathbb R$ as a
real vector space. Its actual real dimension is one:
\[
 d=\dim_{\mathbb R}\mathbb R=1.
\]
There are three agents and three distinct preference orders in
the simultaneous domain. Therefore
\[
 3>2=2^1\cdot1.
\]
Even counting distinct preferences instead of agents does not
rescue the claimed bound.

The profile also has the usual median-voter behavior: alternative
$b$, the peak of the middle voter, defeats $a$ by two votes to one
and defeats $c$ unanimously. Thus the counterexample does not
depend on a failure of median existence.

\textbf{Scope.}
This refutes the stated maximal-agent-count conjunct.
It makes no claim about the other conjectural criteria or the
separate size formula for the full single-peaked domain.
Single-crossing is verified for an explicit voter order, and
single-peakedness for the same fixed alternative axis.

\textbf{Formal correspondence.}
Lean defines actual real locations and negative squared-distance
utilities. An explicit finite rank table represents the three
orders, and the utility/preference theorem proves equivalence
to comparison of the actual utilities for every agent and pair.

The proof verifies strict total preferences and pairwise
distinctness. \texttt{SinglePeaked} is the full increase-to-peak,
decrease-from-peak condition. \texttt{SingleCrossing} requires
every ordered-pair comparison set to be downward or upward closed.
Both properties are proved on the actual finite profile.
The actual \texttt{Module.finrank} of $\mathbb R$ is one,
and the final theorem combines all profile properties with
the strict failure of the displayed count bound.

\textbf{Reproduction.}
Run \texttt{lake build} in \texttt{lean/} with Lean 4.19.0 and
public Mathlib revision
\begin{center}\small\texttt{c44e0c8ee63ca166450922a373c7409c5d26b00b}\end{center}
The project prints axiom audits.
\end{document}
